J'effectue r�guli�rement des interventions br�ves (1 ou 2h) dans des formations � destinations de chercheurs. Je liste ici les formations dans lesquelles ma contribution est plus notable.

\begin{description}
\item[Ecole (JC)2BIM:] cette �cole proposait des cours avanc�s en statistique et algorithmique � de jeunes chercheurs en bioinformatique et biostatistique. Elle s'est tenue au printemps 2018 et sera renouvel�e en 2020 (\url{http://www.gdr-bim.cnrs.fr/?page_id=560}). J'ai contribu� � sa conception et � son organisation en tant que membre du bureau du GdR BIM. Je suis �galement intervenu en tant que formateur.
\item[Formation NGB:] cette formation a �t� organis�e dans le cadre de l'ANR NGB et a accueillie principalement des coll�gue �cologue impliqu� dans cette ANR, ainsi que dans l'ANR EcoNet. Elle portait sur l'analyse des donn�es d'abondance et de r�seau en �cologie et notamment sur les mod�le SBM ({\sl Stochastic blockmodel}) et PLN ({Poisson log-normal}) que je d�cris dans mon rapport de recherche. La demande pour cette formation a �t� tr�s forte et elle devrait �tre renouvel�e dans le cadre de cette ANR ou dans celui du GdR ReSoDiv.
\item[NAIM 2019 :] cette �cole est organis�e conjointement pas la SFdS et l'European Courses in Advances Statistics (ECAS: \url{https://www.sfds.asso.fr/fr/ecas/632-home/}). Elle portera cette ann�e sur Network Analysis, Inference, and Modeling et j'y donnerai une s�rie de cours sur les mod�les graphiques et l'inf�rence de r�seau.
\end{description}

